%! Composition of Functions — Unit 5, Lesson 5.6 — Lesson Plan
\documentclass[10pt]{article}
\usepackage{algebra2-article}
\usepackage{algebra2-boxes}
\usepackage{graphicx}   % -article does NOT load graphicx; needed for thumbnails/screenshots
\graphicspath{{images/}}
\newcommand{\TallMath}[1]{$\displaystyle #1\rule[-1.4em]{0pt}{3.2em}$}
\newcommand{\CourseName}{Algebra 2: Shepherd}
\newcommand{\SchoolYear}{2026--2027}
\newcommand{\MeetingLength}{55 minutes}
\newcommand{\UnitNumberName}{Unit 5: Radical Functions \quad}
\newcommand{\LessonNumberName}{Lesson 5.6: Composition of Functions}

\begin{document}
\begin{center}
    {\Large\bfseries \CourseName: \SchoolYear \\
    {\small\bfseries \UnitNumberName \LessonNumberName}}
\end{center}

\begin{tcolorbox}[colback=forestbg, colframe=forest, boxrule=0.9pt, arc=2mm,
    left=2mm, right=2mm, top=1.5mm, bottom=1.5mm]
    \textbf{Primary Objective:} Students can read the notation $(f\circ g)(x)=f(g(x))$ correctly and
    evaluate a composition \emph{numerically} (from rules, from a table), \emph{graphically} (from
    two graphs), and \emph{algebraically} (by substituting one whole function into the other),
    always working \textbf{inside out} and always recording the intermediate value. They can
    demonstrate that $(f\circ g)$ and $(g\circ f)$ are generally \emph{different functions},
    determine the \textbf{domain} of a composition by checking both stages rather than reading it
    off the simplified rule, and \emph{decompose} a given function into an inner and an outer
    machine. Above all they can explain the lesson's two claims: that order is part of the answer,
    and that a composition \emph{remembers where it has been} --- $\left(\sqrt{x}\,\right)^{2}$
    simplifies to $x$ and still refuses every negative input.
    \\[2pt]
    \textbf{Standards (2023 VA SOL):} \textbf{A2.F.2k} (governing --- determine the composition of
    two functions \emph{algebraically and graphically}; the algebraic strand runs through Notes
    \S4--\S6 and every component, the graphical strand through Notes \S3, Activity Tier~R Part~3, and
    Homework item~4, with the numeric/table strand supporting both). The lesson is the direct
    prerequisite for \textbf{A2.F.2i/j} in Lesson 5.7 --- composition is the \emph{verification} tool
    for inverses ($f(g(x))=g(f(x))=x$), and Tier~A rows (d)/(e) plus Tier~E Part~4 are built to plant
    that. It revisits \textbf{A2.F.1b/c} (5.4's transformations, now revealed as compositions),
    \textbf{A2.F.2a} (domain, computed through two gates), \textbf{A2.EO.2a} (5.1's
    $\sqrt{x^{n}}$ rule, which is why $\sqrt{x^{2}}=|x|$), and \textbf{A2.F.2f} (interpret in
    context --- the homework's mast-climb model).
\end{tcolorbox}

\renewcommand{\arraystretch}{1.4}
\begin{skillbox}[Priority Ideas \& Skills]{goldbox}
\begin{tabularx}{\linewidth}{@{}X|X@{}}
\textbf{Priority Ideas \& Skills} & \textbf{Key Understandings (the ``why'')} \\[2pt]
\hline
\vspace{-6pt}
\begin{itemize}[leftmargin=1.1em, nosep, after=\vspace{2pt}]
  \item Read $(f\circ g)(x)$ correctly: the function \emph{next to the $x$} runs first.
  \item Evaluate a composition from rules, from a table, and from two graphs --- writing the middle number every time.
  \item Compose \textbf{algebraically} by substituting the whole inner function, in parentheses, for every $x$.
  \item Show by counterexample that $(f\circ g)(x)\ne(g\circ f)(x)$ in general.
  \item Find the \textbf{domain} of a composition from the two machines, not from the simplified rule.
  \item Explain why $\left(\sqrt{x}\,\right)^{2}$ has domain $x\ge0$ even though it equals $x$.
  \item \textbf{Decompose} $h(x)$ into an inner and an outer function.
  \item Recognize a pair for which \emph{both} orders return $x$ --- and say what that suggests.
\end{itemize}
&
\vspace{-6pt}
\begin{itemize}[leftmargin=1.1em, nosep, after=\vspace{2pt}]
  \item \emph{Order is part of the answer.} The coupon and the discount are the proof, and the gap is a constant \$$2$ at every price.
  \item Composition is not a new operation --- students did one in Warm-Up 1(d) before the symbol existed.
  \item The wire between the two machines carries exactly one number. A student who stops writing it is a student about to reverse the order.
  \item $\circ$ is not $\cdot$. Two genuinely different misconceptions live in Homework item 5, and they need different corrections.
  \item An input must clear \textbf{two gates}. Simplifying can erase the evidence that a gate was there, but not the gate.
  \item \emph{A composition remembers where it has been.} This is the whole reason 5.7 must restrict the domain of $y=x^{2}$.
  \item Every transformation from 5.4 was a composition --- which is finally \emph{why} inside changes act backwards.
  \item When both orders give back $x$, the two functions are undoing each other. That is tomorrow's definition, discovered today.
\end{itemize}
\end{tabularx}
\end{skillbox}

\begin{skillbox}[{Vocabulary, Concepts, \& Theorems}]{sky}
\renewcommand{\arraystretch}{1.3}
\begin{tabularx}{\linewidth}{@{}l X@{}}
\textbf{Composition of functions} & A function built by running two functions in series: the output of the first becomes the input of the second. \\
\textbf{$(f\circ g)(x)$} & Means $f(g(x))$; read ``$f$ of $g$ of $x$.'' The circle is \emph{not} a multiplication dot. \\
\textbf{Inner function} & The one written next to the $x$ ($g$ above). It runs \textbf{first}. \\
\textbf{Outer function} & The one wrapped around it ($f$ above). It runs \textbf{second}. \\
\textbf{Inside out} & The evaluation order: innermost parentheses first, then work outward. \\
\textbf{The middle number} & The single value passed from the inner machine to the outer one. Write it down every time. \\
\textbf{Domain of a composition} & The inputs that survive both stages: $x$ in the domain of $g$, \emph{and} $g(x)$ in the domain of $f$. \\
\textbf{Gate 1 / Gate 2} & Names for those two tests, used all lesson so students can say \emph{which} stage refused an input. \\
\textbf{Decomposition} & Writing a given $h(x)$ as $f(g(x))$ --- separating the step done first from the step done second. \\
\textbf{Non-commutativity} & $(f\circ g)(x)\ne(g\circ f)(x)$ in general. For two linear functions, equality holds exactly when $ad+b=cb+d$ (Tier E). \\
\end{tabularx}
\end{skillbox}

\begin{fixedskillbox}[Activate Prior Knowledge \& Spiral Review]{forestbg}
    The Warm-Up hands over all three tools and names none of them. \textbf{(1)} Function-notation
    evaluation with $f(x)=2x-3$ and $g(x)=x^{2}+1$ --- and then part (d), which asks students to take
    their answer to $g(-1)$ and feed it into $f$. \emph{That is $(f\circ g)(-1)$, computed before the
    notation exists.} When you introduce the symbol in \S1, point straight back at it: nobody should
    leave thinking composition is a new operation. \textbf{(2)} $h(x)=x^{2}-4x$ evaluated at $3$, at
    $a$, and at $x+1$. Part (c) \emph{is} the algebraic method of \S4 --- substitute the whole
    expression, in parentheses, for every $x$. Police the parentheses here, because the student who
    writes $x^{2}+1-4x+1$ now will write $2x-3^{2}+1$ later. \textbf{(3)} Start at $10$; subtract $4$
    then double ($12$), or double then subtract $4$ ($16$). Get ``order matters'' said out loud by a
    student, on arithmetic nobody can argue with, \emph{before} the jacket comes out at two prices ---
    otherwise half the room will assume a cashier made a mistake. Resolve none of the three; each is
    answered by a specific section of the notes.
\end{fixedskillbox}

\begin{skillbox}[Hook]{forestbg}
    A \$$80$ jacket, two offers running at once: $20\%$ off ($f(x)=0.80x$) and a \$$10$ coupon
    ($g(x)=x-10$). Two cashiers apply \emph{both}, in opposite orders. Walk the money on the board
    and let the class fill the table: coupon first gives $80\rightarrow70\rightarrow\mathbf{56}$;
    discount first gives $80\rightarrow64\rightarrow\mathbf{54}$. Then ask the question that matters:
    \textbf{``Which cashier made the mistake?''} Neither did --- and that is the point. Push once
    more: \emph{where did the missing \$$2$ go?} It is $20\%$ of the coupon; applying the coupon
    first lets the discount eat part of your own savings. (Tier E proves the gap is \$$2$ at
    \emph{every} price, which is worth previewing aloud even if you do not assign it.) Land the
    sentence and keep it on the board all period: \emph{two functions, run one after the other ---
    and the order is part of the answer.} Then name it: today that pairing gets a notation.
\end{skillbox}

\begin{skillbox}[Lesson]{forestbg}
\begin{multicols}{2}
\textbf{1. The notation.} Draw the two machines in series and put $(f\circ g)(x)=f(g(x))$ under
them. Say the rule the way you want it repeated: \emph{the function next to the $x$ goes first}, so
$g$ is the inner machine and $f$ the outer, and \textbf{you work inside out}. Then two warnings that
cost marks all year: $\circ$ is not $\cdot$, and the middle number gets written down every single
time. Close by returning to Warm-Up 1(d) --- they already did $(f\circ g)(-1)=1$.

\textbf{2--3. Numerically, then graphically.} Both sections exist to make ``order matters''
undeniable \emph{before} any algebra. The table trio is the best thirty seconds in the lesson: one
input, $x=0$, three compositions, three answers ($0$, $-2$, $2$). Graphically the move is identical
--- read the inside value on its own graph, carry that number over, read the outside value on the
other. Insist on the language: \emph{an output on one graph becomes an input on the other.}

\textbf{4. Algebraically.} This is Warm-Up item 2 with a function in place of $x+1$. Work
$(f\circ g)(x)=2x^{2}-1$ beside $(g\circ f)(x)=4x^{2}-12x+10$ and then check both at $x=3$ against
the numbers already on the board --- students should see that the rule and the value agree, which is
what makes the algebra trustworthy.

\columnbreak

\textbf{5. Order matters, and the 5.4 payoff.} The anchor: $f(x)=\sqrt{x}$, $g(x)=x-4$. One order
gives $\sqrt{x-4}$, the other $\sqrt{x}-4$ --- two graphs students already read in Lesson 5.4, with
different endpoints \emph{and} different domains. Put both on the board and ask which one
$f(g(x))$ produces \emph{before} computing. Then deliver the payoff: every transformation from 5.4
was a composition. Inside changes are an inner machine getting at the input first, which is finally
\emph{why} they act backwards --- $\sqrt{x-4}$ shifts right because the inner machine hands the root
a number $4$ smaller, so $x$ must be $4$ bigger to compensate. That question has been open since
5.4; close it here.

\textbf{6. Domain --- the heart of the lesson.} Two gates: $x$ must be in $g$'s domain, and $g(x)$
must be in $f$'s. Warm up on $\sqrt{2x-10}$ ($x\ge5$), then run the surprise: with $f(x)=x^{2}$ and
$g(x)=\sqrt{x}$, $\left(\sqrt{x}\,\right)^{2}$ simplifies to $x$ and \emph{still} cannot accept
$-9$, because Gate 1 turned it away before stage 2 ever ran. Beside it, $\sqrt{x^{2}}=|x|$ (5.1),
which accepts everything but is not $x$. \emph{Neither order is honestly ``$x$ for every real
number''} --- do not resolve that; it is tomorrow's lesson.

\textbf{7. Decomposition.} Two minutes, and cuttable. The useful question: what would your hands do
first if you evaluated $h(4)$ on a calculator? That is the inner function.
\end{multicols}
\end{skillbox}

\begin{skillbox}[Active Monitoring]{redbox}
Circulate during the notes practice and the tiers. Watch for and correct:
\begin{itemize}[leftmargin=1.2em, nosep]
  \item \textbf{the signature error --- reversed order.} It is a \emph{reading} error, not a mathematical one, so correct it with the sentence, not with algebra: \emph{which one is next to the $x$?} It is exit-ticket distractor A;
  \item \textbf{no middle number written down.} Treat this as an unfinished problem. A student composing in their head is a student about to reverse the order, and you cannot diagnose an answer with no intermediate step;
  \item \textbf{$\circ$ read as multiplication} --- $x^{2}\cdot(x+3)$. Rare but serious; sit with them and run one composition on the two-box diagram. It is Homework 5 (Amara) and exit-ticket distractor D;
  \item \textbf{dropped parentheses:} $g(2x-3)$ written as $2x-3^{2}+1$. Send them back to Warm-Up 2(c);
  \item \textbf{$(x+1)^{2}=x^{2}+1$} --- the Unit 3 perfect-square error resurfacing inside a composition. It is exit-ticket distractor C;
  \item \textbf{reading the domain off the simplified rule.} The tell is answering ``all reals'' for $\left(\sqrt{x}\,\right)^{2}$ or for Tier A Part 3(c). This is the one that must be fixed \emph{before} 5.7;
  \item \textbf{``a square root means a restricted domain''} --- the wrong generalization, which passes Homework 3(a)/(b) and fails 3(e), where the radicand $x^{2}+2$ can never go negative;
  \item \textbf{$\sqrt{x^{2}}$ written as $x$}, dropping the absolute value. That is 5.1's $\sqrt[n]{a^{n}}$ rule, and today is the only place this week where variables are not assumed non-negative.
\end{itemize}
Cold-call prompts: ``Which machine runs first --- and how do you know by \emph{looking}?'' ``What is
the middle number, and which graph did you read it on?'' ``Before you compute: will the other order
give the same answer? Why not?'' ``This simplified to $x+1$ with no radical in it. Is the domain all
real numbers? What refused, and when?'' ``You got $x$ both ways on that pair. What are those two
functions doing to each other?''
\end{skillbox}

\begin{skillbox}[Group Work \& Differentiation]{redbox}
\textbf{Wired in Series} activity (tiered; mirrors the notes).
\begin{multicols}{3}
\textbf{Tier R --- Remediate.} Part 1 runs $f(x)=x+4$ and $g(x)=3x$ both ways in a table with a
column for \emph{which runs first} and a column for \emph{the middle number} --- $10$ against $18$,
with a required explanation of why tripling last also triples the $+4$. Part 2 is four more with
$2x-1$ and $x^{2}$. Part 3 repeats the notes' two pre-drawn graphs and asks for four reads; rows (a)
and (b) deliberately start from the same input $x=-4$ and finish at $0$ and $2$, so the closing
question --- \emph{why don't they land in the same place?} --- has a concrete answer: different
numbers crossed the wire.

\columnbreak
\textbf{Tier A --- Approaching.} Part 1 is a five-row both-orders table: two ordinary pairs, the
radical pair from \S5, and then rows (d) $x+3$ / $x-3$ and (e) $3x-1$ / $\frac{x+1}{3}$, which come
out to $x$ \emph{both ways}. \textbf{Do not pre-empt this} --- J1 asks students to name what those
pairs are doing to each other and invent a word for it, and you should collect their wordings to
open Lesson 5.7. J2 is the assessed conceptual item: state both domains in row (c) and say
\emph{which stage} refuses in each. J3 diagnoses the order error from a wrong answer. Part 3 is
three domain problems ending in $x^{2}+1$ composed with $\sqrt{x}$, which simplifies to $x+1$ and is
\emph{still} restricted.

\columnbreak
\textbf{Tier E --- Extension.} Four parts. \emph{(1)} The jacket, generalized: $0.8x-8$ against
$0.8x-10$, so the gap is a constant \$$2$ at every price, traced to $0.20\times\$10$. \emph{(2)} The
general linear result --- $f(g(x))=acx+ad+b$, $g(f(x))=acx+cb+d$, so the $x$-terms always agree and
the two commute exactly when $ad+b=cb+d$; tested against Tier A row (d) and the jacket pair.
\emph{(3)} Three machines: $(f\circ g\circ h)(2)=25$, and $2\sqrt{x-5}+1$ taken apart into three
steps, then compared with the 5.4 transformation list --- \emph{a transformation is a composition}.
\emph{(4)} The $\left(\sqrt{x}\,\right)^{2}$ versus $\sqrt{x^{2}}$ table at $x=-4,-1,0,4,9$, ending
on the question that opens 5.7: what has to be done to $y=x^{2}$ first?
\end{multicols}
\end{skillbox}

\begin{skillbox}[Individual Work \& Assessment]{redbox}
\textbf{Exit Ticket} (independent, no notes): with $f(x)=3x-2$ and $g(x)=x^{2}+1$, compute
$(f\circ g)(2)=13$ and $(g\circ f)(2)=17$, then name which function was evaluated first; then
$(f\circ g)(x)=\sqrt{x+7}$ with domain $x\ge-7$, where item~2(c) requires a one-sentence
justification naming \emph{which stage} does the refusing --- that is the conceptual item; then an
\textbf{SOL-style multiple-choice} item, $f(x)=x^{2}$ and $g(x)=x-5$, with options $x^{2}-5$,
$(x-5)^{2}$, $x^{2}-25$, and $x^{3}-5x^{2}$ (answer B). Each distractor is a distinct failure:
\textbf{A} is the order error (it is $(g\circ f)$), \textbf{C} is the Unit 3 perfect-square error
inside a correct composition, \textbf{D} is reading $\circ$ as multiplication. Collect and sort into
``got it / order reversed / domain missing / composition $=$ multiplication.'' \textbf{The pile that
decides tomorrow is the third one} --- students who wrote $\sqrt{x+7}$ and then answered ``all
reals'' are simplifying and reading the domain off the answer, and the entire restricted-domain
argument in 5.7 will be arbitrary to them. Re-run Notes \S6 with that group before starting 5.7.
\end{skillbox}

\begin{skillbox}[Reinforcement \& Extension]{goldbox}
\textbf{Homework:} six numeric compositions from rules (with a follow-up naming the order as the
culprit); six from a \emph{table}, closing with 6(g), the first \emph{backwards} question of the
unit --- for what $x$ is $(p\circ q)(x)=0$? --- which is the seed of 5.7's swap-and-solve; five
algebraic compositions with domains, whose follow-up kills the wrong generalization ``a square root
means a restricted domain'' (row (e) is $\sqrt{x^{2}+2}$, and the radicand can never go negative);
six graphical reads off two pre-drawn graphs, with part (g) noticing that the two orders agree for
$x\ge1$, past the corner of $u$; an error analysis carrying \emph{two different} misconceptions
side by side (Amara multiplies, Ben reverses) settled by evaluating all three at $x=2$; and a
justification item (\textbf{the bridge to 5.7}) explaining why $\left(\sqrt{x}\,\right)^{2}=x$ has
domain $x\ge0$ while $\sqrt{x^{2}}=|x|$ accepts everything. \textbf{Extension:} the horizon model
$d(h)=\sqrt{1.5h}$ from Lesson 5.4 composed with a mast climb $h(t)=6t+6$, giving
$(d\circ h)(t)=3\sqrt{t+1}$ --- whole-number distances at $t=0,3,8,15$ ($3,6,9,12$ miles), a
horizontal shift that \emph{is} the sailor's $6$-foot head start and a vertical stretch that is the
$\sqrt{1.5\cdot6}=3$ hidden in the formula, the quadrupling ($3$ seconds buys the second $3$ miles,
$12$ more buys the next $6$), and a context-domain item rejecting $t=-1$. It closes the loop opened
in 5.4: a real navigation formula is the parent plus a shift and a stretch --- and now students know
\emph{why}, because the shift and the stretch are the two machines they composed.
\textbf{Preview:} Lesson~5.7, \emph{Inverse Functions} --- some of today's pairs gave back $x$ in
\emph{both} orders, and that is not a curiosity, it is the definition. $f(g(x))=g(f(x))=x$ becomes
the test; swap-and-solve builds the partner; the pair reflects over $y=x$ and trades domain for
range; and $\left(\sqrt{x}\,\right)^{2}$ finally explains why $y=x^{2}$ needs a restricted domain ---
the thread first pulled in Lesson 5.0.
\end{skillbox}
\end{document}
